\begin{itemize}
\item \textbf{Kovová vazba:} kladné iontové zbytky spojují společné, volně pohyblivé valenční elektrony.
\item \textbf{Iontová vazba:} přenosem elektronu vznikne kationt a aniont, které se elektrostaticky přitahují; např. NaCl.
\item \textbf{Kovalentní vazba:} atomy sdílejí elektronové páry; při nerovnoměrném sdílení je vazba polární.
\end{itemize}
Rozhodující je \textbf{elektronegativita a její rozdíl $\Delta\chi$}: orientačně $\Delta\chi<0{,}4$ nepolární kovalentní, $0{,}4$--$1{,}7$ polární kovalentní a $>1{,}7$ iontová vazba. Kovová vazba vzniká mezi atomy kovů s nízkou elektronegativitou.
